% Ch7 WS2 -- electron configurations. Block 3.
\documentclass{apchem-worksheet}

\apcunitword{Chapter}
\apcunit{7}
\apcunittitle{Atomic Structure and Periodicity}
\apcdoctitle{Worksheet 2 \textbullet\ Electron Configurations}
\apcsubtitle{Zumdahl \S7.11 \textbullet\ build the ion first, then configure it}

\begin{document}
\wsheader[Write configurations from memory using the periodic table's block
structure, then check. For ions, adjust the electron count before you write
anything.]

\problem[]{Write full electron configurations:}
\begin{parts}
  \item \ce{N}: \answerblank[5cm]{$1s^2 2s^2 2p^3$}
  \item \ce{Al}: \answerblank[5cm]{$1s^2 2s^2 2p^6 3s^2 3p^1$}
  \item \ce{Ar}: \answerblank[5cm]{$1s^2 2s^2 2p^6 3s^2 3p^6$}
\end{parts}

\problem[]{Write noble-gas shorthand configurations:}
\begin{parts}
  \item \ce{Se}: \answerblank[4.6cm]{$[\ce{Ar}]4s^2 3d^{10} 4p^4$}
  \item \ce{Sr}: \answerblank[4.6cm]{$[\ce{Kr}]5s^2$}
  \item \ce{Mn}: \answerblank[4.6cm]{$[\ce{Ar}]4s^2 3d^5$}
  \item \ce{I}: \answerblank[4.6cm]{$[\ce{Kr}]5s^2 4d^{10} 5p^5$}
\end{parts}

\problem[]{Ions --- adjust the electron count first:}
\begin{parts}
  \item \ce{Ca^2+}: \answerblank[4cm]{$[\ce{Ar}]$}
  \item \ce{P^3-}: \answerblank[4cm]{$[\ce{Ar}]$}
  \item \ce{Fe^2+}: \answerblank[4cm]{$[\ce{Ar}]3d^6$}
  \item \ce{Fe^3+}: \answerblank[4cm]{$[\ce{Ar}]3d^5$}
  \item \ce{Zn^2+}: \answerblank[4cm]{$[\ce{Ar}]3d^{10}$}
\end{parts}

\problem[]{Transition-metal cations lose 4$s$ electrons before 3$d$, even
though 4$s$ filled first. State the rule and apply it to explain why
\ce{Fe^2+} is $[\ce{Ar}]3d^6$ rather than $[\ce{Ar}]4s^2 3d^4$.
\answerlines[3]{Once the 3$d$ orbitals are occupied they drop below 4$s$ in
energy, so the 4$s$ electrons are now the outermost and least tightly held
--- they leave first. Removing two electrons from Fe ($[\ce{Ar}]4s^2 3d^6$)
therefore empties 4$s$ and leaves $3d^6$ untouched.}}

\problem[]{Unpaired electrons. Give the number in the ground state:}
\begin{parts}
  \item \ce{C}: \answerblank[1.8cm]{2}
  \item \ce{O}: \answerblank[1.8cm]{2}
  \item \ce{Fe^3+}: \answerblank[1.8cm]{5}
  \item \ce{Zn}: \answerblank[1.8cm]{0}
\end{parts}

\problem[]{Classify each as ground state, excited state, or impossible ---
and identify the species where one exists:}
\begin{parts}
  \item $1s^2 2s^2 2p^6 3s^1$:
        \answerblank[5cm]{ground --- Na}
  \item $1s^2 2s^2 2p^6 3s^2 3p^5 4s^1$ (18 e$^-$):
        \answerblank[5cm]{excited Ar}
  \item $1s^2 2s^2 2p^6 3s^3$:
        \answerblank[5cm]{impossible ($s$ max 2)}
  \item $1s^2 2s^1 2p^1$ (4 e$^-$):
        \answerblank[5cm]{excited Be}
\end{parts}

\problem[]{The two exceptions:}
\begin{parts}
  \item Predicted vs actual configuration of \ce{Cr}:
        \answerblank[6cm]{$[\ce{Ar}]4s^2 3d^4$ vs $[\ce{Ar}]4s^1 3d^5$}
  \item Predicted vs actual configuration of \ce{Cu}:
        \answerblank[6cm]{$[\ce{Ar}]4s^2 3d^9$ vs $[\ce{Ar}]4s^1 3d^{10}$}
  \item Give the reason both anomalies occur.
        \answerlines[2]{Half-filled and completely filled $d$ subshells are
        slightly more stable, and the 4$s$ and 3$d$ energies are close
        enough that promoting one electron is favourable.}
\end{parts}

\problem[]{Name three species isoelectronic with \ce{Ne} and state their
shared configuration.
\answerlines[2]{Any three of \ce{F-}, \ce{O^2-}, \ce{N^3-}, \ce{Na+},
\ce{Mg^2+}, \ce{Al^3+} --- all are $1s^2 2s^2 2p^6$ with 10 electrons.}}

\begin{spiralstrip}{Chapter 7 \textbullet\ blocks 1--2}
  \item Energy of a \SI{400}{\nano\meter} photon:
        \answerblank[3.2cm]{\SI{4.97e-19}{\joule}}
  \item Below threshold frequency, brighter light does what?
        \answerblank[2.8cm]{nothing}
  \item Why is 3$s$ below 3$p$? \answerblank[4.1cm]{greater penetration}
  \item Orbitals in the $n=3$ shell: \answerblank[1.6cm]{9}
  \item $Z_{\text{eff}}$ stands for: \answerblank[4.8cm]{effective nuclear
        charge}
\end{spiralstrip}

\end{document}
